\FigureLayoutDeclare{img/2005/jiangxi_liberal/q21_fig1.png}{8.0cm}{43bp 65bp 72bp 20bp}
\chapter{2005年江西卷（文）}
\section{单选题}
\begin{problem}
设集合 \(I = \{x \mid |x| < 3, x \in \Z\}\)，\(A = \{1, 2\}\)，\(B = \{-2, -1, 2\}\)，则 \(A \cup (\complement_{I} B) =\)（\quad）

\choices
    {\(\{1\}\)}
    {\(\{1, 2\}\)}
    {\(\{2\}\)}
    {\(\{0, 1, 2\}\)}
\end{problem}

\begin{answer}
D
\end{answer}

\begin{problem}
已知 \(\tan \frac{\alpha}{2} = 3\)，则 \(\cos \alpha =\)（\quad）

\choices
    {\(\frac{1}{5}\)}
    {\(-\frac{4}{5}\)}
    {\(\frac{4}{15}\)}
    {\(-\frac{3}{5}\)}
\end{problem}

\begin{answer}
B
\end{answer}

\begin{problem}
\((\sqrt{x} + \sqrt[3]{x})^{12}\) 的展开式中，含 \(x\) 的正整数次幂的项共有（\quad）

\choices
    {4 项}
    {3 项}
    {2 项}
    {1 项}
\end{problem}

\begin{answer}
B
\end{answer}

\begin{problem}
函数 \(f(x) = \frac{1}{\log_{2}(-x^{2} + 4x - 3)}\) 的定义域为（\quad）

\choices
    {\((1, 2) \cup (2, 3)\)}
    {\((-\infty, 1) \cup (3, +\infty)\)}
    {\((1, 3)\)}
    {\([1, 3]\)}
\end{problem}

\begin{answer}
A
\end{answer}

\begin{problem}
设函数 \(f(x) = \sin 3x + |\sin 3x|\)，则 \(f(x)\) 为（\quad）

\choices
    {周期函数，最小正周期为 \(\frac{\pi}{3}\)}
    {周期函数，最小正周期为 \(\frac{2\pi}{3}\)}
    {周期函数，最小正周期为 \(2\pi\)}
    {非周期函数}
\end{problem}

\begin{answer}
B
\end{answer}

\begin{problem}
已知向量 \(\bs{a} = (1, 2)\)，\(\bs{b} = (-2, -4)\)，\(|\bs{c}| = \sqrt{5}\)，若 \((\bs{a} + \bs{b}) \cdot \bs{c} = \frac{5}{2}\)，则 \(\bs{a}\) 与 \(\bs{c}\) 的夹角为（\quad）

\choices
    {\(30^{\circ}\)}
    {\(60^{\circ}\)}
    {\(120^{\circ}\)}
    {\(150^{\circ}\)}
\end{problem}

\begin{answer}
C
\end{answer}

\begin{problem}
将 9 个（含甲，乙）平均分成三组，甲，乙分在同一组，则不同分组方法的种数为（\quad）

\choices
    {70}
    {140}
    {280}
    {840}
\end{problem}

\begin{answer}
A
\end{answer}

\begin{problem}
在 \(\triangle ABC\) 中，设命题 \(p: \frac{a}{\sin B} = \frac{b}{\sin C} = \frac{c}{\sin A}\)，命题 \(q: \triangle ABC\) 是等边三角形，那么命题 \(p\) 是命题 \(q\) 的（\quad）

\choices
    {充分不必要条件}
    {必要不充分条件}
    {充分必要条件}
    {既不充分又不必要条件}
\end{problem}

\begin{answer}
C
\end{answer}

\begin{problem}
矩形 \(ABCD\) 中，\(AB = 4\)，\(BC = 3\)，沿 \(AC\) 将矩形 \(ABCD\) 折成一个直二面角 \(B - AC - D\)，则四面体 \(ABCD\) 的外接球的体积为（\quad）

\choices
    {\(\frac{125}{12}\pi\)}
    {\(\frac{125}{9}\pi\)}
    {\(\frac{125}{6}\pi\)}
    {\(\frac{125}{3}\pi\)}
\end{problem}

\begin{answer}
C
\end{answer}

\begin{problem}
已知实数 \(a\)，\(b\) 满足等式 \(\left(\frac{1}{2}\right)^a = \left(\frac{1}{3}\right)^b\)，下列五个关系式：

\begin{circlelist}
\item \(0 < b < a\)；
\item \(a < b < 0\)；
\item \(0 < a < b\)；
\item \(b < a < 0\)；
\item \(a = b\).
\end{circlelist}

其中不可能成立的关系式（\quad）

\choices
    {1 个}
    {2 个}
    {3 个}
    {4 个}
\end{problem}

\begin{answer}
B
\end{answer}

\begin{problem}
在 \(\triangle OAB\) 中，\(O\) 为坐标原点，\(A(1, \cos \theta)\)，\(B(\sin \theta, 1)\)，\(\theta \in \left(0, \frac{\pi}{2}\right]\)，则 \(\triangle OAB\) 的面积达到最大值时，\(\theta =\)（\quad）

\choices
    {\(\frac{\pi}{6}\)}
    {\(\frac{\pi}{4}\)}
    {\(\frac{\pi}{3}\)}
    {\(\frac{\pi}{2}\)}
\end{problem}

\begin{answer}
D
\end{answer}

\begin{problem}
为了解某校高三学生的视力情况，随机地抽查了该校 \(100\) 名高三学生的视力情况，得到频率分布直方图. 由于不慎将部分数据丢失，但知道前 \(4\) 组的频数成等比数列，后 \(6\) 组的频数成等差数列，设最大频率为 \(a\)，视力在 \(4.6\) 到 \(5.0\) 之间的学生数为 \(b\)，则 \(a\)，\(b\) 的值分别为（\quad）

\bitmapfigure[max width=0.74\linewidth]{img/2005/jiangxi_liberal/q12_fig1.png}

\choices
    {0.27, 78}
    {0.27, 83}
    {2.7, 78}
    {2.7, 83}
\end{problem}

\begin{answer}
A
\end{answer}

\section{填空题}
\begin{problem}
若函数 \(f(x) = \log_a\left(x + \sqrt{x^2 + 2a^2}\right)\) 是奇函数，则 \(a =\)\fillinblank{}.
\end{problem}

\begin{answer}
\(\frac{\sqrt{2}}{2}\)
\end{answer}

\begin{problem}
设实数 \(x\)，\(y\) 满足 \(\begin{cases}
x - y - 2 \leqslant 0,\\
x + 2y - 4 > 0,\\
2y - 3 \leqslant 0,
\end{cases}\) 则 \(\frac{y}{x}\) 的最大值是\fillinblank{}.
\end{problem}

\begin{answer}
不存在
\end{answer}

\begin{solution}
由 \(x-y-2\leqslant0\) 和 \(x+2y-4>0\)，分别得
\(x\leqslant y+2\) 和 \(x>4-2y\). 要使这样的 \(x\) 存在，必须有
\[
4-2y<y+2,
\]
从而 \(y>\frac{2}{3}\). 再结合 \(2y-3\leqslant0\)，得
\(\frac{2}{3}<y\leqslant\frac{3}{2}\). 因而 \(4-2y>0\)，且
\[
0<\frac{y}{x}<\frac{y}{4-2y}\leqslant\frac{3}{2}.
\]
但取任意 \(0<\varepsilon\leqslant\frac{5}{2}\)，令
\(y=\frac{3}{2}\)，\(x=1+\varepsilon\)，则三式均成立，且
\[
\frac{y}{x}=\frac{3}{2(1+\varepsilon)}\longrightarrow\frac{3}{2}
\quad(\varepsilon\to0^+)
\]
所以 \(\frac{y}{x}\) 的上确界为 \(\frac{3}{2}\)，但上确界不能取到，故不存在最大值.
\end{solution}

\begin{problem}
如图，在三棱锥 \(P - ABC\) 中，\(PA = PB = PC = BC\)，且 \(\angle BAC = \frac{\pi}{2}\)，则 \(PA\) 与底面 \(ABC\) 所成角为\fillinblank{}.

\bitmapfigure[max width=0.42\linewidth]{img/2005/jiangxi_liberal/q15_fig1.png}
\end{problem}

\begin{answer}
\(\frac{\pi}{3}\)
\end{answer}

\begin{problem}
以下四个关于圆锥曲线的命题中：

\begin{circlelist}
\item 设 \(A\)，\(B\) 为两个定点，\(k\) 为非零常数，\(\left| \overrightarrow{PA} \right| - \left| \overrightarrow{PB} \right| = k\)，则动点 \(P\) 的轨迹为双曲线；
\item 设圆 \(C\) 上一定点 \(A\)，作圆的动点弦 \(AB\)，\(O\) 为坐标原点，若 \(\overrightarrow{OP} = \frac{1}{2} \left(\overrightarrow{OA} + \overrightarrow{OB}\right)\)，则动点 \(P\) 的轨迹为椭圆；
\item 方程 \(2x^2 - 5x + 2 = 0\) 的两根可分别作为椭圆和双曲线的离心率；
    \item 双曲线 \(\frac{x^2}{25} - \frac{y^2}{9} = 1\) 与椭圆 \(\frac{x^2}{35} + y^2 = 1\) 有相同的焦点.
\end{circlelist}

其中真命题的序号为\fillinblank{}.（写出所有真命题的序号）
\end{problem}

\begin{answer}
\(\circled{3}\circled{4}\)
\end{answer}

\section{解答题}
\begin{problem}
已知函数 \(f(x) = \frac{x^2}{ax + b}\)（\(a\)，\(b\) 为常数），且方程 \(f(x) - x + 12 = 0\) 有两个实根为 \(x_1 = 3\)，\(x_2 = 4\).

\begin{enumerate}
\item 求函数 \(f(x)\) 的解析式；
\item 设 \(k > 1\)，解关于 \(x\) 的不等式：\(f(x) < \frac{(k + 1)x - k}{2 - x}\).
\end{enumerate}
\end{problem}

\begin{solution}
\begin{enumerate}
\item 将 \(x=3\) 和 \(x=4\) 代入方程，分别得
\(\frac{9}{3a+b}-3+12=0\)，\(\frac{16}{4a+b}-4+12=0\).
因此 \(3a+b=-1\)，\(4a+b=-2\)，解得 \(a=-1\)，\(b=2\)，所以 \(f(x)=\dfrac{x^2}{2-x}\)，且 \(x\ne2\).
\item 由函数定义域知 \(x\ne2\)，原不等式等价于
\[
\frac{x^2}{2-x}-\frac{(k+1)x-k}{2-x}<0
\quad\Longleftrightarrow\quad
\frac{(x-1)(x-k)}{2-x}<0.
\]
当 \(1<k<2\) 时，按 \(1\)，\(k\)，\(2\) 分点讨论，得 \(x\in(1,k)\cup(2,+\infty)\)；当 \(k=2\) 时，得 \(x\in(1,2)\cup(2,+\infty)\)；当 \(k>2\) 时，得 \(x\in(1,2)\cup(k,+\infty)\).
\end{enumerate}
\end{solution}

\begin{problem}
已知向量 \(\bs{a} = \left(2\cos \frac{x}{2}, \tan \left(\frac{x}{2} + \frac{\pi}{4}\right)\right)\)，\(\bs{b} = \left(\sqrt{2}\sin \left(\frac{x}{2} + \frac{\pi}{4}\right), \tan \left(\frac{x}{2} - \frac{\pi}{4}\right)\right)\)，令 \(f(x) = \bs{a} \cdot \bs{b}\). 求函数 \(f(x)\) 的最大值，最小正周期，并写出 \(f(x)\) 在 \([0, \pi]\) 上的单调区间.
\end{problem}

\begin{solution}
令 \(t=\frac{x}{2}\). 两个切线函数都有定义时，
\[
2\cos t\cdot\sqrt{2}\sin\left(t+\frac{\pi}{4}\right)=1+\sin x+\cos x
\]
且
\[
\tan\left(t+\frac{\pi}{4}\right)\tan\left(t-\frac{\pi}{4}\right)=-1
\]
所以在定义域
\[
\mathcal D=\left\{x\in\R\mid x\ne\frac{\pi}{2}+n\pi,\ n\in\Z\right\}
\]
上，\(f(x)=\sin x+\cos x=\sqrt{2}\sin\left(x+\frac{\pi}{4}\right)\). 因此函数的最大值为 \(\sqrt{2}\). 定义域中被排除的点相邻间隔为 \(\pi\)，故函数的正周期必须是 \(\pi\) 的整数倍；又 \(f(x+\pi)=-f(x)\)，而 \(f(x+2\pi)=f(x)\)，所以最小正周期为 \(2\pi\). 在 \([0,\pi]\) 上，原函数在 \(x=\frac{\pi}{2}\) 处无定义，且 \(f'(x)=\cos x-\sin x\)，故在 \([0,\frac{\pi}{4}]\) 上单调递增，在 \([\frac{\pi}{4},\frac{\pi}{2})\) 和 \((\frac{\pi}{2},\pi]\) 上单调递减.
\end{solution}

\begin{problem}
\(A\)，\(B\) 两位同学各有五张卡片，现以投掷均匀硬币的形式进行游戏，当出现正面朝上时 \(A\) 赢得 \(B\) 一张卡片，否则 \(B\) 赢得 \(A\) 一张卡片，如果某人已赢得所有卡片，则游戏终止. 求掷硬币的次数不大于 7 次时游戏终止的概率.
\end{problem}

\begin{solution}
设正面出现的次数与反面出现的次数之差为 \(D_n\). 每投掷一次硬币，\(D_n\) 增加 \(1\) 或减少 \(1\)，游戏在 \(D_n=5\) 或 \(D_n=-5\) 时终止.

由于 \(D_n\) 与 \(n\) 具有相同的奇偶性，游戏在 \(7\) 次以内终止的次数只能是 \(5\) 或 \(7\). \(5\) 次终止时，必须连续出现 \(5\) 次正面或连续出现 \(5\) 次反面，所以
\[
P(T=5)=\frac{2}{2^5}=\frac{1}{16}
\]
在第 \(7\) 次终止时，以正面终止为例，前 \(6\) 次中必须有 \(5\) 次正面和 \(1\) 次反面，且反面必须出现在前 \(5\) 个位置之一，否则会在第 \(5\) 次已经终止. 因此共有 \(5\) 种；以反面终止同理有 \(5\) 种，从而
\[
P(T=7)=\frac{10}{2^7}=\frac{5}{64}
\]
所以掷硬币的次数不大于 \(7\) 次时游戏终止的概率为
\[
P(T\le7)=\frac{1}{16}+\frac{5}{64}=\frac{9}{64}
\]
\end{solution}

\begin{problem}
如图，在长方体 \(ABCD - A_{1}B_{1}C_{1}D_{1}\) 中，\(AD = AA_{1} = 1\)，\(AB = 2\)，点 \(E\) 在棱 \(AB\) 上移动.

\begin{enumerate}
\item 证明：\(D_{1}E \perp A_{1}D\)；
\item 当 \(E\) 为 \(AB\) 的中点时，求点 \(E\) 到面 \(ACD_{1}\) 的距离；
\item \(AE\) 等于何值时，二面角 \(D_{1} - EC - D\) 的大小为 \(\frac{\pi}{4}\).
\end{enumerate}

\bitmapfigure[max width=0.55\linewidth]{img/2005/jiangxi_liberal/q20_fig1.png}
\end{problem}

\begin{solution}
建立空间直角坐标系，使 \(D\) 为原点，\(DA\)、\(DC\)、\(DD_{1}\) 分别为 \(x\)、\(y\)、\(z\) 轴正方向. 设 \(AE=t\)，则 \(0\leqslant t\leqslant2\)，各点坐标为
\[
D(0,0,0),\quad A(1,0,0),\quad C(0,2,0),\quad D_{1}(0,0,1),\quad A_{1}(1,0,1),\quad E(1,t,0).
\]

\begin{enumerate}
\item
\(\overrightarrow{D_{1}E}=(1,t,-1)\)，\(\overrightarrow{A_{1}D}=(-1,0,-1)\)，所以
\[
\overrightarrow{D_{1}E}\cdot\overrightarrow{A_{1}D}=-1+1=0.
\]
故 \(D_{1}E\perp A_{1}D\).

\item 当 \(E\) 为 \(AB\) 的中点时，\(E(1,1,0)\). 平面 \(ACD_{1}\) 的一个法向量为
\[
\overrightarrow{AC}\times\overrightarrow{AD_{1}}=(2,1,2)
\]
故其方程为
\[
2x+y+2z-2=0
\]
于是点 \(E\) 到该平面的距离为
\[
\frac{|2\times1+1+2\times0-2|}{\sqrt{2^{2}+1^{2}+2^{2}}}=\frac{1}{3}
\]

\item 平面 \(D_{1}EC\) 的一个法向量为
\[
\overrightarrow{D_{1}E}\times\overrightarrow{D_{1}C}=(2-t,1,2)
\]
而平面 \(DEC\) 即底面 \(ABCD\)，可取法向量 \((0,0,1)\). 设二面角 \(D_{1}-EC-D\) 为 \(\theta\)，则
\[
\cos\theta=\frac{2}{\sqrt{(2-t)^{2}+1+4}}.
\]
当 \(\theta=\frac{\pi}{4}\) 时，
\[
\frac{2}{\sqrt{(2-t)^{2}+5}}=\frac{\sqrt{2}}{2}
\quad\Longrightarrow\quad
(2-t)^{2}=3.
\]
由于 \(0\leqslant t\leqslant2\)，所以 \(2-t\geqslant0\)，从而
\[
t=2-\sqrt{3}
\]
即 \(AE=2-\sqrt{3}\).
\end{enumerate}
\end{solution}

\begin{problem}
如图，\(M\) 是抛物线 \(y^2 = x\) 上的一点，动弦 \(ME\)，\(MF\) 分别交 \(x\) 轴于 \(A\)，\(B\) 两点，且 \(MA = MB\).

\begin{enumerate}
\item 若 \(M\) 为定点，证明：直线 \(EF\) 的斜率为定值；
\item 若 \(M\) 为动点，且 \(\angle EMF = 90^{\circ}\)，求 \(\triangle EMF\) 的重心 \(G\) 的轨迹方程.
\end{enumerate}

\bitmapfigure[max width=0.50\linewidth]{img/2005/jiangxi_liberal/q21_fig1.png}
\end{problem}

\FloatBarrier

\begin{solution}
\begin{enumerate}
\item 设 \(M=(u^2,u)\)，\(E=(v^2,v)\)，\(F=(w^2,w)\). 抛物线上参数为 \(r\)，\(s\) 的两点连线斜率为 \(\dfrac{1}{r+s}\)，且该直线与 \(x\) 轴交于 \((-rs,0)\). 因此 \(A=(-uv,0)\)，\(B=(-uw,0)\). 由于 \(A\)，\(B\) 是两个不同的交点，故 \(u\ne0\)，从而
\[
MA^2=u^2\left((u+v)^2+1\right),
\qquad
MB^2=u^2\left((u+w)^2+1\right).
\]
由 \(MA=MB\)，并且 \(E\ne F\)，得 \((u+v)^2=(u+w)^2\)，故 \(v+w=-2u\). 于是直线 \(EF\) 的斜率为
\[
\frac{1}{v+w}=-\frac{1}{2u},
\]
只由定点 \(M\) 决定.
\item 由 \(v+w=-2u\)，令 \(s=u+v\)，则 \(u+w=-s\). 若 \(s=0\)，则 \(v=w=-u\)，与 \(E\ne F\) 矛盾，故 \(s\ne0\). 两弦 \(ME\)，\(MF\) 的斜率分别为 \(\dfrac{1}{s}\)，\(-\dfrac{1}{s}\). 由 \(\angle EMF=90^{\circ}\)，得
\[
\frac{1}{s}\cdot\left(-\frac{1}{s}\right)=-1,
\qquad s^2=1
\]
所以 \(s=\pm1\)，并且 \(v=s-u\)，\(w=-s-u\). 设重心 \(G=(x_G,y_G)\)，则
\[
x_G=\frac{u^2+(s-u)^2+(-s-u)^2}{3}=u^2+\frac{2}{3},
\qquad
y_G=\frac{u+(s-u)+(-s-u)}{3}=-\frac{u}{3}.
\]
消去 \(u\)，得 \(x_G=9y_G^2+\frac{2}{3}\)，即
\[
y_G^2=\frac{1}{9}x_G-\frac{2}{27}.
\]
由于 \(A\)，\(B\) 为两个不同的交点，\(u\ne0\)，故轨迹为 \(y^2=\dfrac{1}{9}x-\dfrac{2}{27}\)，其中 \(x>\dfrac{2}{3}\).
\end{enumerate}
\end{solution}

\begin{problem}
已知数列 \(\{a_n\}\) 的前 \(n\) 项和 \(S_n\) 满足 \(S_n - S_{n - 2} = 3\left(-\frac{1}{2}\right)^{n - 1}\)（\(n \geqslant 3\)），且 \(S_1 = 1\)，\(S_2 = -\frac{3}{2}\). 求数列 \(\{a_n\}\) 的通项公式.
\end{problem}

\begin{solution}
由 \(a_1=S_1=1\)，\(a_2=S_2-S_1=-\frac{5}{2}\)，且对 \(n\ge3\)，
\[
a_{n-1}+a_n=S_n-S_{n-2}=3\left(-\frac{1}{2}\right)^{n-1}
\]
因此
\[
a_n=3\left(-\frac{1}{2}\right)^{n-1}-a_{n-1}
\]
下面用归纳法求 \(a_n\). 当 \(n=1\) 和 \(n=2\) 时，分别有 \(a_1=4-3\)，\(a_2=-4+3\left(\frac{1}{2}\right)\)，结论成立. 若结论对 \(n-1\) 成立，当 \(n\) 为奇数时，\(n-1\) 为偶数，代入递推式得
\[
a_n=3\left(\frac{1}{2}\right)^{n-1}+4-3\left(\frac{1}{2}\right)^{n-2}=4-3\left(\frac{1}{2}\right)^{n-1}
\]
当 \(n\) 为偶数时，\(n-1\) 为奇数，代入递推式得
\[
a_n=-3\left(\frac{1}{2}\right)^{n-1}-4+3\left(\frac{1}{2}\right)^{n-2}=-4+3\left(\frac{1}{2}\right)^{n-1}
\]
所以
\[
a_n=
\begin{cases}
4-3\left(\dfrac{1}{2}\right)^{n-1}, & n\text{ 为奇数},\\
-4+3\left(\dfrac{1}{2}\right)^{n-1}, & n\text{ 为偶数}.
\end{cases}
\]
\end{solution}
